\chapter{序論}
\chapoverview{本章では，本研究の背景，問題意識，目的，研究課題，貢献を述べる．まず，特許文書を用いた技術機会探索の重要性を整理する．次に，既存の特許文書埋め込み研究では，課題文脈と解決手段文脈が混在しやすいという問題を示す．その上で，本研究がU-Netの医療画像解析から製造欠陥検出への展開を対象に，課題・解決手段分離型のクロスドメイン技術候補抽出を検討する研究であることを述べる．}

\section{背景}
特許文書は，企業や研究機関の研究開発活動を反映する技術情報である．特許文書には，発明の目的，課題，解決手段，効果，対象物，出願人，分類コードなどが記録されている．そのため，特許情報は技術動向分析，競合分析，研究開発テーマ探索，技術機会分析に広く用いられてきた\cite{griliches1990patent, oecd2009patent}．

異なる技術分野の間であっても，取り組む課題が類似する場合がある．たとえば，医療画像解析と製造欠陥検出は，対象物は人体画像と製造物画像で異なる．しかし，画像中の異常領域，欠陥領域，対象領域を高精度に検出または分割するという点では共通する．このような分野間では，一方の分野で発展した技術手法が，他方の分野でも後年に利用される可能性がある．

近年，BERTやSentence-BERTなどの文書埋め込み手法により，特許文書をベクトル空間上に写像し，文書間の意味的類似性を分析する研究が進んでいる\cite{reimers2019sentencebert, bekamiri2021patentsberta}．文書埋め込みを用いることで，従来のキーワード一致や分類コードだけでは捉えにくい意味的な近さを扱える．このため，特許文書埋め込みは，類似特許検索，特許分類，技術ランドスケープ分析，技術機会探索に応用されている．

\section{問題意識}
既存の特許文書埋め込み研究では，特許全文を1つのベクトルとして扱い，意味的に近い特許を探索することが多い．しかし，特許全文ベクトルでは，発明が「何を解決しようとしているのか」と「どのように解決しているのか」が混在する．そのため，異分野応用候補を抽出する際に，候補が課題の類似によって選ばれたのか，解決手段の類似によって選ばれたのか，あるいは対象物や用語の一致によって選ばれたのかが説明しにくい．

クロスドメイン技術候補を抽出するには，単に意味的に近い特許を探すだけでは不十分である．対象分野と参照分野が似た課題を扱っているか，その課題に対して参照分野ではどのような解決手段が用いられているか，さらにその解決手段が対象分野では低出現または未使用であるかを分けて確認する必要がある．

本研究では，この問題に対して，特許文書を課題文脈と解決手段文脈に分けて扱う．課題文脈は，従来技術の問題点，解決したい課題，目的，必要性を表す．解決手段文脈は，その課題を解決するための構成，方法，材料，アルゴリズム，処理手順を表す．この2つを分けることで，異分野間の関係を「課題の近さ」と「解決手段の出現差分」として整理できる．

\section{研究目的}
本研究の目的は，特許文書から課題文脈と解決手段文脈を分離して抽出し，課題が近い異分野間において，対象分野では低出現または未使用である他分野の技術手法を，専門家が確認すべきクロスドメイン技術候補として抽出できるかを検討することである．

本研究では，医療画像解析を参照分野，製造欠陥検出を対象分野とし，U-Netを主ケースとして扱う．U-Netは生物医学画像セグメンテーション向けに提案された技術手法であり\cite{ronneberger2015unet}，後年には製造欠陥検出や表面欠陥セグメンテーションにも応用されている．そのため，既知のクロスドメイン応用事例を後ろ向きに確認する対象として適している．

\section{研究課題}
本研究では，以下の3つの研究課題を設定する．

\begin{description}[leftmargin=3.0em]
  \item[RQ1] 特許文書を課題文脈と解決手段文脈に分離することで，異分野技術候補の抽出理由を説明可能にできるか．
  \item[RQ2] 2015--2018年の過去期間において，医療画像解析で先行して出現するU-Netを，製造欠陥検出分野における低出現候補として抽出できるか．
  \item[RQ3] 過去期間に候補として抽出されたU-Netは，2019--2024年の将来評価期間において，製造欠陥検出分野での出現・増加と対応するか．
\end{description}

\section{本研究の貢献}
本研究の貢献は以下の3点である．

第1に，特許文書を課題文脈と解決手段文脈に分け，クロスドメイン技術候補抽出を「課題の近さ」と「解決手段の低出現性」の関係として整理する点である．

第2に，U-Netを対象として，医療画像解析と製造欠陥検出の特許データを比較し，過去期間の手法ギャップと将来期間の出現を用いた後ろ向き評価設計を示す点である．

第3に，Publication numberベースとFamily IDベースの両方で傾向を確認し，代表特許の目視確認を組み合わせることで，単なる語の一致ではなく，技術的手段としての利用を確認する点である．

\section{本論文の構成}
本論文の構成は以下の通りである．第2章では，特許データを用いた技術分析，特許文書埋め込み，技術機会分析，クロスドメイン技術探索に関する関連研究を整理する．第3章では，課題・解決手段分離に基づく候補抽出枠組みを示す．第4章では，使用データと実験設計を説明する．第5章では，U-Net予備分析の結果を示す．第6章では，結果の解釈，限界，今後の展開を議論する．第7章では，本研究の結論を述べる．
